\begin{abstractCN}
投资者即使持有相同的期末财富，也可能因为此前经历了上涨回落或下跌反弹，而以不同的参考点评价下一次投资。累积前景理论（CPT）能够描述这种参考依赖的风险判断，但现有的一阶策略优化通常固定参考点；在连续投资中，策略改变财富路径，财富路径又改变参考点，使下一期的CPT目标随市场和账户经历共同移动。本文研究这一情形下的在线投资组合学习，提出动态参考点累积前景策略梯度方法（DRCPT-PG）。

方法以股票因子驱动的Dirichlet策略生成满足权重约束的随机投资组合，沿每条财富轨迹逐日更新参考点，并用期末相对结果评价策略。针对每一期给定状态后的CPT目标，本文结合中心化留一估计与随机层级修正，构造严格无偏的梯度估计量，证明其具有有限期望轨迹成本和有限方差，并给出渐近正态性及固定维度下的最优轨迹预算速率。跨期采用归一化投影更新跟踪移动目标；当目标总变差为次线性、估计与模拟误差的平均影响逐渐消失，平均平方投影残差趋于零，动态局部遗憾呈次线性增长。

实验依次考察参考点的路径依赖、梯度估计和在线跟踪。在连续缓变的半合成市场中，持续更新策略的累积平方投影残差为$0.684$，冻结策略为$1.427$。按时间顺序进行的沪深300真实市场回放表明，在五种策略梯度方法中，DRCPT-PG取得最高的平均Sharpe比率，以及最低的平均最大回撤、换手和交易成本。真实投资者的样本外持股决策进一步显示，非对称动态参考点对减持行为的预测比静态参考点更准确。由此，参考点既记录投资经历对盈亏判断的影响，也成为策略更新时需要随账户路径持续跟踪的状态。

\keywordsCN{累积前景理论；动态参考点；策略梯度；投资组合优化；动态局部遗憾}
\end{abstractCN}

\begin{abstractEN}
Cumulative prospect theory (CPT) models portfolio choice relative to a reference point, yet most first-order CPT policy methods keep that reference fixed even when realized wealth should update it.  We study online portfolio optimization with an endogenous reference state, where policy decisions reshape wealth, wealth updates the reference, and the resulting CPT objective moves over time.  We introduce Dynamic-Reference Cumulative Prospect Policy Gradient (DRCPT-PG), which combines an exactly unbiased estimator of the regularized frozen-target CPT gradient with finite expected trajectory cost and variance, and a normalized projected update for tracking the moving first-order target.  The estimator supports studentized asymptotic inference and attains the fixed-dimensional minimax trajectory-budget rate, while the online analysis controls the projected residual and dynamic local regret under sublinear objective variation.  In controlled experiments, continued updating reduces cumulative squared projected residual from $1.427$ to $0.684$ relative to a frozen closed-loop control.  In a chronological 300-stock A-share replay, DRCPT-PG achieves the highest mean Sharpe ratio and the lowest mean maximum drawdown, turnover, and transaction cost among five policy-gradient methods; held-out investor data further show incremental predictive gains from adaptive and asymmetric reference representations.  These results connect adaptive reference formation to a statistically controlled first-order learning framework for path-dependent online portfolio optimization.
\end{abstractEN}
